% This file was generated automatically.
% Source: judge.txt

\begin{prompttextbox}
\begingroup
\ttfamily
\small
\sloppy
\obeyspaces
An LLM agent attempted to answer a question relating to some time series data. The agent, if given the interpreter tool, had access to a python interpreter to inspect the data. Here is the full conversation, including thinking traces (surrounded by <think></think> tags) and tool calls (surrounded by <tool code></tool code> tags):\par
\par\smallskip
=== Agent's conversation ===\par
\par\smallskip
\{data[conv]\}\par
\par\smallskip
=== Conversation end ===\par
\par\smallskip
\par\smallskip
The correct answer was \{data[correct]\}.\par
\par\smallskip
Output a json matching the template below to answer the questions about the agent's conversation. Always explain your decisions under the "X\_expl" field in 1-3 sentences.\par
\endgroup

\par\smallskip
\begin{lstlisting}[style=promptjsonstyle]
{
// "core_reason_for_answer" should be one of the following options:
//  - "code", if the agent answered mainly based on code output.
//  - "raw data", if the agent answered mainly based on "eyeballing" the raw data.
//  - "question format", if the agent answered mainly based on strategizing about the
// exam format.
//  - "other", if none of the options above works.
"core_reason_for_answer":
  enum["code", "raw data", "question format", "other"],
"core_reason_for_answer_expl": string,

// "core_problem_solving_strategy" should be one of the following options, based on the
// primary strategy of the agent:
// - "statistical test", if the primary strategy was to perform a formal hypothesis
// test or compute a test statistic and p-value (examples: Augmented Dickey--Fuller,
// KPSS, Granger causality, t-test, Wilcoxon test, chi-square test, permutation test).
// - "spectral analysis", if the primary strategy was to analyze frequency content of
// the series (examples: FFT, periodogram, Welch's method, Lomb--Scargle, bandpower,
// spectral peak detection).
// - "curve fitting", if the primary strategy was to fit a parametric model to data
// (examples: least-squares polynomial/regression fitting, AR/ARMA/ARIMA/SARIMA
// parameter estimation, nonlinear curve fits, sklearn or statsmodels model.fit usage,
// model-based forecasting).
// - "windowed/rolling statistics", if the primary strategy was to compute statistics
// over sliding or rolling windows, without transforming into the frequency domain
// (examples: rolling mean/variance, rolling correlations or cross-correlations,
// rolling regression, windowed autocorrelation, rolling volatility).
// - "simple arithmetic", if the primary strategy consisted of direct, non-model-based
// summary calculations without statistical inference or signal analysis (examples:
// global mean/median/variance, counts, ratios, totals, expanding statistics, basic
// differencing).
// - "other", if none of the above fits.
"core_strategy":
  enum["statistical test", "spectral analysis", "curve fitting", "windowed/rolling statistics", "simple arithmetic", "other"],
"core_strategy_expl": string,

"methodological_problems": {
  // "conceptual_misunderstanding" should be true if the agent demonstrates a fundamental
  // misunderstanding of key concepts relevant to the question.
  "conceptual_misunderstanding": bool,
  "conceptual_misunderstanding_expl": optional[string],

  // "wrong_core_problem_solving_strategy" should be true if the strategy chosen by the
  // agent ("core_problem_solving_strategy") was not suitable for answering the question.
  // For example, if the model decided to use curve fitting, when using a statistical
  // test would have been better.
  "wrong_core_problem_solving_strategy": bool,
  "wrong_core_problem_solving_strategy_expl":
  optional[string],

  // "wrong_method_within_strategy" should be true if the strategy chosen by the agent
  // ("core_problem_solving_strategy") was suitable for answering the question, but
  // failed due to the specific method chosen within the strategy. For example, the model
  // selected an incorrect statistical test, heuristic, threshold, or algorithmic
  // variant.
  "wrong_method_within_strategy": bool,
  "wrong_method_within_strategy_expl": optional[string],

  // "unsupported_assumption" should be true if for the method that the agent chose, it
  // made an unjustified assumption about the data or the options. For example, assuming
  // stationarity without evidence.
  "unsupported_assumption": bool,
  "unsupported_assumption_expl": optional[string],

  // "implementation_errors" should be true if there were coding or algorithmic mistakes
  // that affect the validity of the result. Code execution errors are that caused the
  // code to fail are not included in this category.
  "implementation_errors": bool,
  "implementation_errors_expl": optional[string],

  // "incorrect_result_interpretation" should be true if the agent incorrectly interprets
  // the result of any of his computations. For example, the model wrongly conclude that
  // the difference is statistically significant.
  "incorrect_result_interpretation": bool,
  "incorrect_result_interpretation_expl": optional[string],

  // "insufficient_evidence_guess" should be true if the agent answered whilst not having
  // gathered sufficiently convincing evidence.
  "insufficient_evidence_guess": bool,
  "insufficient_evidence_guess_expl": optional[string]
},

// This section is primarily meant for agents that had access to a code interpreter
// tool.
"code_problems": {
  // "code_result" should be one of the following options, indicating the overall success
  // of the code the model produced (if any):
  // - "success", if the agent managed to obtain the data it intended. If the code
  // initially failed due to errors, but the agent managed to fix the errors, it still
  // counts as a success.
  // - "partial failure", if the agent failed to obtain some of the desired data, but
  // managed to get at enough results to produce an answer. For example if the code
  // failed after printing sufficient evidence.
  // - "complete failure", if the agent failed to obtain the results and gave up or ran
  // out of turns.
  // - "no code", if the agent didn't produce any code.
  "code_result":
  enum["success", "partial failure", "complete failure", "no code"],
  "code_result_expl": optional[string],

  // "tool_usage_trouble" should be true if the agent had some trouble following
  // instructions to use the tool correctly. This mainly includes cases where the agent
  // generated an improper coding tool call or had trouble calling the tool properly.
  // Regular errors and exceptions that caused the code to fail are not included in this
  // category.
  "tool_usage_trouble": bool,
  "tool_usage_trouble_expl": optional[string],

  // "other" should be true if there was some other issue with the code that does not
  // fall into any of the other categories.
  "other": bool,
  "other_expl": optional[string]
},

"other_problems": {
  // "reasoning_answer_mismatch" should be true if the agent settled on one answer in the
  // thinking trace and then answered differently.
  "reasoning_answer_mismatch": bool,
  "reasoning_answer_mismatch_expl": optional[string],

  // "reasoning_tool_usage_mismatch" should be true if the agent claims it will use the
  // coding tool (such as saying "let's compute", or crafting a code solution in the
  // thinking trace) but then goes straight to answering.
  "reasoning_tool_usage_mismatch": bool,
  "reasoning_tool_usage_mismatch_expl": optional[string],

  // "hallucinated_values_in_reasoning" should be true if the agent introduces numerical
  // values that do not appear in the data or the code output. Reasonable rounding does
  // not count.
  "hallucinated_values_in_reasoning": bool,
  "hallucinated_values_in_reasoning_expl": optional[string]
}
},
\end{lstlisting}

\end{prompttextbox}
